% sections/04_axis_specifications.tex
\section{Axis Specifications}\label{sec:axis-specifications}

This section specifies each of the six thematic axes as implemented in the
\ISI v\,1.0 backend.  For every axis, the specification covers: (i)~the
economic concept measured, (ii)~the data source and product/flow scope,
(iii)~the channel decomposition and aggregation variant, (iv)~the
fallback-basis logic, and (v)~boundary conditions.

% =========================================================================
\subsection{Axis~1: Financial Sovereignty}\label{sec:axis:finance}

\paragraph{Concept.}
Concentration of a country's cross-border financial exposure---banking claims
and portfolio debt---across counterparty countries.

\paragraph{Sources.}
\begin{itemize}
  \item \textbf{Channel~A}: BIS Locational Banking Statistics (Table~A6.2,
    by counterparty country).  Quarterly frequency; the most recent quarter
    available for the vintage year is used (typically Q4).
  \item \textbf{Channel~B}: IMF Coordinated Portfolio Investment Survey
    (CPIS), debt securities.  Annual frequency.
\end{itemize}

\paragraph{Channel decomposition.}
\begin{itemize}
  \item \textbf{Channel~A (banking claims):}  Shares of inward banking
    claims on country~\(i\) by creditor country.
    \(s_{i,j}^{(A)} = \text{claims}_{j \rightarrow i} \;/\;
    \text{claims}_{\cdot \rightarrow i}\).
  \item \textbf{Channel~B (portfolio debt):}  Shares of country~\(i\)'s
    outward portfolio-debt holdings by debtor country.
    \(s_{i,j}^{(B)} = \text{debt}_{i \rightarrow j} \;/\;
    \text{debt}_{i \rightarrow \cdot}\).
\end{itemize}

\paragraph{Cross-channel aggregation.}
Arithmetic mean (\cref{eq:dual-channel}).

\paragraph{Fallback basis.}
When both channels are available: \texttt{BOTH}.  When only BIS data exist:
\texttt{A\_ONLY}.  When only CPIS data exist: \texttt{B\_ONLY}.  When
neither source reports data for a country: \texttt{NONE}; axis score is null.

\paragraph{Boundary conditions.}
Countries absent from the BIS reporting population receive
\(C_i^{(A)} = \text{null}\).  Countries absent from the CPIS receive
\(C_i^{(B)} = \text{null}\).  If both channels are null the axis score is
null and flagged \texttt{W-ZER}.

% =========================================================================
\subsection{Axis~2: Energy Dependency}\label{sec:axis:energy}

\paragraph{Concept.}
Concentration of a country's fossil-fuel energy imports across partner
countries.

\paragraph{Source.}
Eurostat energy tables (\texttt{nrg\_ti\_*}).  Annual frequency.  Scope:
natural gas, crude oil, and solid fossil fuels.

\paragraph{Aggregation variant.}
Axis~2 does \emph{not} use a two-channel decomposition.  Instead, a
per-fuel \HHI is computed for each fossil-fuel category, and the axis
score is the arithmetic average across fuels with non-zero imports
(\cref{eq:fuel-avg}):
%
\[
  M_i^{(2)}
  \;=\;
  \frac{1}{|\mathcal{F}_i|}
  \sum_{f \in \mathcal{F}_i} C_i^{(f)}
\]
%
where \(\mathcal{F}_i \subseteq \{\text{gas},\, \text{oil},\,
\text{solid fossil}\}\) is the set of fuel types for which
country~\(i\) has non-zero imports.

\paragraph{Boundary conditions.}
Countries with zero fossil-fuel imports across all three categories
receive axis score null and are flagged \texttt{W-ZER}.  Countries with
imports in only one or two categories are scored over the available
fuels (denominator adjusts automatically).

% =========================================================================
\subsection{Axis~3: Technology and Semiconductor Dependency}%
\label{sec:axis:technology}

\paragraph{Concept.}
Concentration of a country's imports of semiconductor devices and
electronic components across partner countries.

\paragraph{Source.}
Eurostat Comext (dataset \texttt{ds-045409}).  Product scope: HS~8541
(semiconductor devices) and HS~8542 (electronic integrated circuits).

\paragraph{Channel decomposition.}
\begin{itemize}
  \item \textbf{Channel~A (aggregate):}  \HHI of partner-country
    shares across combined HS~8541\,+\,8542 imports.
  \item \textbf{Channel~B (category-weighted):}  Per-category \HHI
    (HS~8541 and HS~8542 separately), combined via the value-weighted
    formula (\cref{eq:weighted-b}).
\end{itemize}

\paragraph{Cross-channel aggregation.}
Arithmetic mean (\cref{eq:dual-channel}).

\paragraph{Fallback basis.}
\texttt{BOTH} when aggregate and per-category data are available (the
normal case for all EU-27 states).  \texttt{A\_ONLY} / \texttt{B\_ONLY} /
\texttt{NONE} as for Axis~1.

\paragraph{Boundary conditions.}
Countries with zero recorded semiconductor imports receive axis
score null.

% =========================================================================
\subsection{Axis~4: Defence Industrial Dependency}\label{sec:axis:defence}

\paragraph{Concept.}
Concentration of a country's conventional arms imports across supplier
countries.

\paragraph{Source.}
SIPRI Arms Transfers Database (ATD).  Unit of measurement: SIPRI Trend
Indicator Values (TIV), which provide a common metric for comparing
transfers of different weapons systems.

\paragraph{Temporal window.}
Defence transfers are lumpy.  To mitigate year-to-year volatility, a
\textbf{six-year} rolling window is applied: share vectors are computed
from the sum of TIV flows over years \(2019\)--\(2024\) for v\,1.0.

\paragraph{Channel decomposition.}
\begin{itemize}
  \item \textbf{Channel~A (aggregate):}  \HHI computed over partner-country
    shares of total TIV imports across the six-year window.
  \item \textbf{Channel~B (capability-block weighted):}  Per-capability-block
    \HHI (e.g., combat aircraft, armoured vehicles, naval vessels, missiles,
    air-defence systems), combined via the value-weighted formula
    (\cref{eq:weighted-b}).  Capability blocks follow the SIPRI category
    taxonomy.
\end{itemize}

\paragraph{Cross-channel aggregation.}
Arithmetic mean (\cref{eq:dual-channel}).

\paragraph{Warning codes.}
\texttt{W-DEF} (sparse transfer record: raised when country~\(i\) records
fewer than three transfer events in the six-year window, yielding an \HHI
that may reflect sampling noise rather than structural dependence).

\paragraph{Boundary conditions.}
Countries with zero recorded arms imports on both channels are assigned
axis score null and excluded from the composite via the coverage rule
(\cref{sec:agg:missing}).

% =========================================================================
\subsection{Axis~5: Critical Inputs / Raw Materials}%
\label{sec:axis:critical-inputs}

\paragraph{Concept.}
Concentration of a country's imports of critical raw materials across
partner countries.

\paragraph{Source.}
Eurostat Comext at CN8 level.  The material scope covers a curated list
of critical raw materials aligned with the EU Critical Raw Materials Act
(e.g., rare earths, lithium, cobalt, manganese, tungsten, gallium,
germanium, titanium, and platinum-group metals).

\paragraph{Channel decomposition.}
\begin{itemize}
  \item \textbf{Channel~A (aggregate):}  \HHI computed over partner-country
    shares of the combined import value of all critical materials.
  \item \textbf{Channel~B (material-group weighted):}  Per-material-group
    \HHI, combined via the value-weighted formula (\cref{eq:weighted-b}).
    Material groups are defined in the methodology registry.
\end{itemize}

\paragraph{Cross-channel aggregation.}
Arithmetic mean (\cref{eq:dual-channel}).

\paragraph{Fallback basis.}
As for Axis~1 (\texttt{BOTH} / \texttt{A\_ONLY} / \texttt{B\_ONLY} /
\texttt{NONE}).

\paragraph{Boundary conditions.}
Countries with zero recorded critical-material imports receive axis
score null.

% =========================================================================
\subsection{Axis~6: Logistics and Freight Dependency}%
\label{sec:axis:logistics}

\paragraph{Concept.}
Concentration of a country's freight transport linkages across partner
countries and transport modes.

\paragraph{Source.}
Eurostat transport statistics.  Scope: maritime, road, rail, and air
freight flows (tonnage).

\paragraph{Channel decomposition.}
\begin{itemize}
  \item \textbf{Channel~A (mode \HHI):}  For each transport mode, the \HHI
    is computed over partner-country shares of tonnage.  The mode-level
    scores are then combined as a tonnage-weighted average across modes.
  \item \textbf{Channel~B (partner \HHI per mode):}  For each mode, the
    partner-level \HHI is computed and tonnage-weighted across modes,
    analogous to the category-weighted variant (\cref{eq:weighted-b}).
\end{itemize}

\paragraph{Cross-channel aggregation.}
Arithmetic mean (\cref{eq:dual-channel}).

\paragraph{Boundary conditions.}
Landlocked countries (AT, CZ, HU, LU, SK) have no maritime freight data;
the maritime mode is excluded from their tonnage-weighted average
(denominator adjusts automatically).  Countries with complete absence of
transport data receive axis score null.

% =========================================================================
\subsection{Axis Summary}\label{sec:axis:summary}

\Cref{tab:axis-summary} provides a compact reference.

\begin{table}[htbp]
  \centering\small
  \caption{Axis registry summary.}\label{tab:axis-summary}
  \begin{tabularx}{\textwidth}{@{}c l l l X@{}}
    \toprule
    \textbf{Axis} & \textbf{Domain} & \textbf{Ch.\,A} & \textbf{Ch.\,B}
      & \textbf{Aggregation variant} \\
    \midrule
    1 & Financial sovereignty  & BIS inward claims   & IMF CPIS debt  & Arithmetic mean \\
    2 & Energy dependency      & \multicolumn{2}{l}{Per-fuel \HHI (no channels)} & Fuel average \\
    3 & Technology / semiconductor & Aggregate \HHI  & Category-wtd   & Arithmetic mean \\
    4 & Defence industrial     & Aggregate \HHI      & Capability-wtd & Arithmetic mean (6-yr) \\
    5 & Critical inputs        & Aggregate \HHI      & Material-wtd   & Arithmetic mean \\
    6 & Logistics / freight    & Mode \HHI           & Partner \HHI   & Tonnage-weighted \\
    \bottomrule
  \end{tabularx}
\end{table}